\section{Exact parent laws and return coefficients}
\label{app:parent}

For $c=\delta=d\in\{1,2\}$ and $\beta=\gamma=1$, RNA immigrants form a
Poisson process of rate $d$. Each RNA has an independent exponential
lifetime of rate $d$, translates at rate one while alive, and gives proteins
with independent exponential lifetimes of rate one. This labelled
construction provides a positive stationary law on the infinite lattice.

The number of alive RNAs at observation is Poisson with mean one.
An alive RNA has age distributed as $\operatorname{Exp}(d)$, so the
mean number of its surviving proteins is $v=1-e^{-\mathrm{age}}$ with
density $d(1-v)^{d-1}$ on $[0,1]$. Its mark distribution is
\begin{equation}
 q_d(k)=\frac d{k!}\int_0^1e^{-v}v^k(1-v)^{d-1}\,\mathrm dv.
 \label{eq:qmark}
\end{equation}
Dead RNAs supply independent compound-Poisson protein counts with
L\'evy coefficients and total intensity
\begin{align}
 \lambda_d(k)&=\frac d{k!}\int_0^1e^{-v}v^{k-1}(1-v)^d\,\mathrm dv,
                  \qquad k\geq1,\nonumber\\
 \Lambda_d&=d\int_0^1\frac{(1-e^{-v})(1-v)^d}{v}\,\mathrm dv .
\end{align}
The joint stationary generating function is consequently
\begin{equation}
 G_d(x,y)=
 \exp\left[xq_d(y)-1+\sum_{k\geq1}\lambda_d(k)(y^k-1)\right],
 \label{eq:jointpgf}
\end{equation}
where $q_d(y)=\sum_kq_d(k)y^k$. For $z=y-1$,
\begin{align}
 q_d(1+z)&=d!\sum_{k\geq0}\frac{z^k}{(d+k)!},\nonumber\\
 \log G_{\mathrm{dead},d}(1+z)
 &=d\,d!\sum_{k\geq1}\frac{z^k}{k(d+k)!}.
\end{align}
These expressions obey $zq_d'=d+(z-d)q_d$ and
$z(\log G_{\mathrm{dead},d})'=d(q_d-1)$.
Substitution in the full stationary forward equation
\begin{equation}
 d(x-1)G+[d(1-x)+x(y-1)]G_x+(1-y)G_y=0
\end{equation}
checks the defining reaction law, including catalytic translation.
The known constitutive moments are
\begin{equation}
 \E M=\E P=1,\quad
 \Var(M)=1,\quad
 \Cov(M,P)=\frac1{d+1},\quad
 \Var(P)=1+\frac1{d+1}.
\end{equation}
They reproduce Eq.~\eqref{eq:baseline}; they are inputs, not the
noise-tradeoff contribution.

Let $\pi_{d,p}$ be the protein marginal. It is compound Poisson with
coefficients $\ell_k=q_d(k)+\lambda_d(k)$ for $k\geq1$:
\begin{equation}
 \pi_{d,0}=e^{-1+q_d(0)-\Lambda_d},\qquad
 \pi_{d,p}=\frac1p\sum_{k=1}^p k\ell_k\pi_{d,p-k}.
 \label{eq:marginal}
\end{equation}
No higher probability coefficient enters this recurrence for a lower
coefficient. All probabilities are positive. The alive-RNA Poisson
Palm identity gives $\E[M\ind_{\{P=p\}}]=(\pi_d*q_d)_p$.
Alive RNAs with zero protein marks form an independent Poisson
population of mean $q_d(0)\leq1$. Each other alive RNA accounts for at
least one protein, proving $\E[M\mid P=p]\leq p+1$.

\subsection{Extra-RNA response}

Add one fresh independently marked RNA at time zero to a stationary
parent, and let $K_t$ be its surviving protein count at time $t$. Linearity
of the parent channels gives
\begin{equation}
 D_d(t;h)=\E[h(P+K_t)-h(P)],
\end{equation}
where $P$ is stationary and independent of the added family.
Conditionally on the extra RNA lifetime, $K_t$ is Poisson with mean
between zero and one. For $k\geq1$, direct mark integration gives
\begin{equation}
 J_k(s):=\int_0^\infty e^{-st}\E[(K_t)_k]\dd
 =\frac{k!}{(s+k)(s+d)^{\overline k}},\qquad J_0(s)=\frac1s.
\end{equation}
Here $(x)_k$ is a falling factorial and
$(s+d)^{\overline k}=\prod_{j=0}^{k-1}(s+d+j)$.
Differentiation gives
\begin{equation}
 -J_k'(s)=J_k(s)
 \left[\frac1{s+k}+\sum_{j=0}^{k-1}\frac1{s+d+j}\right].
\end{equation}
Define measures
$\nu_s(k)=\int e^{-st}\Pr(K_t=k)\dd$ and
$\nu_{s,t}(k)=\int te^{-st}\Pr(K_t=k)\dd$.
Their coefficient expansion is
\begin{equation}
 \nu_s(k)=\sum_{l\geq0}
 \frac{(-1)^l\binom{k+l}{k}}
 {(s+k+l)(s+d)^{\overline{k+l}}},
 \label{eq:nu}
\end{equation}
with each term for $\nu_{s,t}(k)$ additionally multiplied by the bracket
in $-J_{k+l}'(s)/J_{k+l}(s)$. Even and odd Taylor sums for the
conditional Poisson factor $e^{-v_t}$ enclose these integrals because
$0\leq v_t\leq1$.

The coefficient sequences of the return functionals are
\begin{align}
 A_p&=(\pi_d*q_d)_p-\pi_{d,p},&
 H_p&=(p-1)\pi_{d,p},\nonumber\\
 (D_s)_p&=(\pi_d*\nu_s)_p-\pi_{d,p}/s,&
 (D_{s,t})_p&=(\pi_d*\nu_{s,t})_p-\pi_{d,p}/s^2.
 \label{eq:coefficients}
\end{align}
All have zero total sum. In Eq.~\eqref{eq:ranktwo},
$D_{g,d}=D_1$;
$D_{r,1}=D_1-\tfrac12D_2$ and
$D_{r,2}=D_1+D_{1,t}$.

For $k\geq1$,
$q_d(k)\geq e^{-1}d!/(d+k)!$. Comparing this with $J_k(s)/k!$,
and using $e<3$, shows
$\nu_1(k),\nu_2(k),\nu_{1,t}(k)\leq(3/2)q_d(k)$.
For the time-weighted case the extra harmonic bracket is at most
$(k+1)/2$. Splitting off count zero in the convolutions, using total
measure mass at most one and the conditional RNA bound, yields
\begin{equation}
 |\ell_p|\leq C_\ell(p+1)\pi_{d,p},\qquad
 \begin{array}{c|cccc}
 \ell&A&H&D_g&D_r\\ \hline
 d=1&1&1&5/2&15/4\\
 d=2&1&1&5/2&5
 \end{array}
 \label{eq:palmconstants}
\end{equation}
The tabulated entries are $C_\ell$. Together with
Eq.~\eqref{eq:dual}, these inequalities control each functional on the
entire response space.

\section{Infinite-state numerical enclosures}
\label{app:certificate}

\subsection{Probabilities, moments and resolvent residuals}

The stationary cluster integrals are enclosed by the $160$th and
$161$st alternating Taylor sums, with outward fixed-point arithmetic
at $384$ bits. Positive convolution recurrences evaluate the joint
law~\eqref{eq:jointpgf} and marginal~\eqref{eq:marginal}.
The joint calculation uses
\begin{equation}
 \pi_d(m,p)=\frac{e^{-1}}{m!}
              [y^p]q_d(y)^mG_{\mathrm{dead},d}(y).
\end{equation}
No finite stationary solve or artificial normalization supplies its
probabilities.

For $f_{r,s}=\E[(M)_r(P)_s]$, direct generator testing gives
\begin{equation}
 (rd+s)f_{r,s}
 =rd\,f_{r-1,s}
       +s\big[f_{r+1,s-1}+r f_{r,s-1}\big],
 \qquad f_{0,0}=1.
 \label{eq:factorialrec}
\end{equation}
The recurrence is triangular at each total degree, proceeding from
smaller $s$. Exact mixed factorial moments through degree $104$
bound the required tails. For example, factorial Markov bounds at
$R+1$ and $S+1$, together with
\begin{align}
 M(M)_k&=(M)_{k+1}+k(M)_k,\nonumber\\
 M^2(M)_k&=(M)_{k+2}+(2k+1)(M)_{k+1}+k^2(M)_k,
\end{align}
provide upper bounds $\tau_j$ on
$\E[M^j\ind_{\{M>R\ {\rm or}\ P>S\}}]$ for $j=0,1,2$.
The $P$-tail terms use the corresponding mixed moments.
Positive bounds are minimized over available factorial orders
$k\leq100$ without moment closure.

Construct approximate resolvent vectors on $R=48,S=128$ using the
capped backward generator. Quantize them as integers over $2^{192}$,
then extend them to the full lattice by coordinate clipping. This cap
defines a witness construction, not the stationary process. For each
source category, the source itself equals its clipped extension outside
the rectangle. If $\eta$ is the maximum exact interior residual and
$V$ the approximant sup norm, then the extended residual obeys
\begin{equation}
 |r(m,p)|\leq\eta+
   2V[dR+S+(d+1)m]\ind_{\{m>R\ {\rm or}\ p>S\}}.
\end{equation}
Missing capped death terms and excess translation account for this
bound. Clipping in $p$ prevents an unbounded $p$ coefficient.
With $A=dR+S$ and $B=d+1$,
\begin{align}
 A_{\rm tail}&=A^2\tau_0+2AB\tau_1+B^2\tau_2,\nonumber\\
 \norm{r}_2^2&\leq
 \sum_{\rm inside}\pi_d^+(m,p)r(m,p)^2
           +2\eta^2\tau_0+8V^2A_{\rm tail}.
 \label{eq:l2residual}
\end{align}
Only exact nonnegative sums and positive tail bounds enter this
enclosure.

Conditional Jensen and invariance make the full Markov semigroup an
$L^2(\pi_d)$ contraction. Hence
$\norm{(s-L_d)^{-1}}_{2\to2}\leq1/s$. For the exact resolvent $v_s$
and its clipped approximant $\widehat v_s$,
\begin{equation}
 \norm{v_s-\widehat v_s}_2\leq\norm{r_s}_2/s.
\end{equation}
For $w=(1-L_d)^{-2}b$, whose approximant uses $\widehat v_1$ as source,
\begin{equation}
 \norm{w-\widehat w}_2\leq\norm{r_1}_2+\norm{r_w}_2.
\end{equation}
Localization and finite polynomial moments justify these residual
identities for the clipped vectors.

Use disjoint protein categories $1,\ldots,n-1$ and $\{P\geq n\}$,
omitting zero modulo constants. Scale each by an exact power of two
$s_k$ so that $1/4<s_k^2\pi_k^+\leq1$.
Testing a residual against a centered scaled category costs at most
one in $L^2$. The injection-matrix resolvent errors are therefore
controlled by $2\norm{r_1}_2+\norm{r_w}_2$ for $U_1$, and
$\norm{r_1}_2+\norm{r_2}_2/4$ for $U_2$.
The second expression includes the $-1/2$ kernel coefficient and
the $1/2$ resolvent norm. Probability integration and centering
errors are propagated separately before symmetrization.

Return matrices use the one-dimensional coefficient
sequences~\eqref{eq:coefficients}. The full threshold coefficient is
the negative prefix sum through $p=n-1$, since each sequence has
zero total sum. No derivative at an RNA boundary is used.

\subsection{Exact block positivity}

Let $G_2,G_1$ be the enclosed injection matrices on
$\mathcal S_{64}$ and $\mathsf B$ the matrix for $V_1$ or $V_2$ in
the corresponding $(h_1,h_2)$ ordering. The required block is
\begin{equation}
 K=\begin{pmatrix}
      (3/4)G_2&\mathsf B\\
      \mathsf B^{\mathsf T}&(3/4)G_1
    \end{pmatrix}.
 \label{eq:block}
\end{equation}
A floating Cholesky inverse proposes a congruence $T$, then its
entries are quantized as integers over $2^{96}$. The central block
$K_0$ is stored as integers over $2^{192}$; its rounding is added to
a common true-entry error bound $e$. If $\ell_i$ is the absolute
row sum of the rational $T$, each transformed entry error is at most
$e\ell_i\ell_j$. The exact sufficient margins are
\begin{equation}
 (TK_0T^{\mathsf T})_{ii}
 -\sum_{j\ne i}|(TK_0T^{\mathsf T})_{ij}|
 -e\ell_i\sum_j\ell_j>0.
 \label{eq:margin}
\end{equation}
The two $128$-dimensional blocks have all these rational margins
positive; their minimum margins both exceed $0.9948$.
Symmetric strict diagonal dominance gives positivity of the true
transformed block. It also forces $T$ to have full rank, proving
positivity of the true original block. Rescaling the two input
vectors in Eq.~\eqref{eq:block} then gives Eq.~\eqref{eq:finitebound}.
No eigenvalue rounded to a positive decimal decides this test.

\subsection{Riesz residuals and full return norms}

Finite Riesz vectors at $n=32,48$ are proposed by a floating solve,
then quantized to integers over $2^{128}$. The $n=64$ Gram enclosures
are transformed by an exact basis map to the smaller spaces; the
new threshold is the sum of all old threshold categories.
Source intervals and Gram errors enclose every scaled residual
in Eq.~\eqref{eq:head}.

For completeness, the polarized injection operator has coefficient
\begin{equation}
 (U_dv)_p=\frac{\pi_{d,p}}2\int_0^\infty r_{\rm rec}(t)
 \left[
 \E[v(P_t)-\E v\mid P_0=p]
 +\E[v(P_0)-\E v\mid P_t=p]
 \right]\dd ,
\end{equation}
where $r_{\rm rec}=r_2$ for $d=1$, and $r_1$ for $d=2$.
Both conditional expectations have absolute value at most $V$.
The kernel masses are $3/4$ and $2$, giving the bound used
in Eq.~\eqref{eq:tail} without assuming reversibility.

Full moments for the tail subtraction can also be reconstructed from
the exact protein factorial cumulants
\begin{equation}
 \kappa_k=\frac{(k-1)!\,d!}{(d+k-1)!}.
\end{equation}
The Bell recurrence gives factorial moments, and Stirling conversion
gives $\E[(P+1)^j]$. Subtracting the lower probability prefix yields
the upper $T_j$ for $j=4,5,6$. This is an independent moment
representation of the same parent law.

All residual products and tail subtractions are rational. Square-root
upper bounds are obtained with integer square roots and checked by
squaring. Orthogonal projection contracts the dual residual norm:
Eq.~\eqref{eq:rieszerror} bounds
$\ell-\ell\mathcal P_{d,n}$ even though the finite spaces need not
exhaust $\mathcal H_d$. Equations~\eqref{eq:ranktwo}
and~\eqref{eq:rankerror} then give the full return bounds in
Table~\ref{tab:certificate}. The stored rational data and
standard-library replay code specify the complete finite arithmetic.

\section{Finite-feedback covariance remainder}
\label{app:remainder}

This appendix proves Eq.~\eqref{eq:remainder} without assuming a
differentiable stationary distribution. All response functions here
retain their original sup bound one, even when constants are removed
for the response-space analysis.

\subsection{A stationary immigration-envelope coupling}

For gene $i$, use common RNA immigration rate $c_i-\epsilon g$ and
an additional envelope immigration process of rate $2\epsilon g$,
where $(c_1,c_2)=(1,2)$. Each arrival is independently marked with
RNA lifetime, translation events and protein lifetimes. A constitutive
parent $Y_i$ accepts each envelope arrival with probability $1/2$.
The regulated process $X_i$ accepts it with probability
\begin{equation}
 \frac{f_i-(c_i-\epsilon g)}{2\epsilon g}\in[0,1],
\end{equation}
using the same uniform acceptance mark. Acceptance depends on the
current regulated state, but the descendant marks of the new family
are independent of that decision. Let $E_i$ be the process retaining
every envelope family. Then, pathwise,
\begin{equation}
 |X_{M,i}-Y_{M,i}|\leq M_{E,i},\qquad
 |X_{P,i}-Y_{P,i}|\leq P_{E,i}.
 \label{eq:envelope}
\end{equation}
The parents remain independent across genes. The two envelopes are
also independent, and $Y_i$ is independent of $E_j$ for $j\ne i$.
The parent and its own envelope share the parent's half-envelope;
their covariance must be retained. Stationary joint limits follow
from constructions started in the past, uniform dominating
polynomial moments and uniqueness of the marginal invariant laws.
No monotone or associated feedback process is required.

For $W=M+P$ in a constitutive gene with transcription $c$ and
$\delta=d\geq1$, the exact moments give
\begin{equation}
 \E W=\frac{2c}{d},\qquad
 \Var(W)=\frac cd\left(2+\frac3{d+1}\right).
\end{equation}
Thus $\E W_{Y,i}=2$, $\E W_{E,i}\leq4\epsilon g$ and
$\Cov(W_{Y,i},W_{E,i})\leq7\epsilon g/2$.
For either $i$,
\begin{equation}
 \E[(1+W_{Y,i})(W_{E,1}+W_{E,2})]
 \leq \frac{55}{2}\epsilon g.
 \label{eq:weightedmoment}
\end{equation}
The constant is $3(4+4)+7/2$; it includes the parent's
dependence on its own envelope.

\subsection{Product-parent Poisson solution}

Write $\eta_j=\E_0h_j(P_j)$ and $\bar z_i=1+1/\delta_i$.
For $i\ne j$, let
\begin{align}
 q_i&=(Z_i-\bar z_i)(h_j(P_j)-\eta_j),\nonumber\\
 u_i(t,m_i,p_i)&=r_i(t)(m_i-1)+e^{-t}(p_i-1),\nonumber\\
 v_j(t,m_j,p_j)&=\mathcal T_j(t)h_j(m_j,p_j)-\eta_j,\nonumber\\
 Rq_i&=\int_0^\infty u_i(t,m_i,p_i)v_j(t,m_j,p_j)\dd .
\end{align}
The affine parent mean evolution and independence give $L_0Rq_i=-q_i$.
Here $|v_j|\leq2$ and $r_i(t)\leq(1+t)e^{-t}$, so $Rq_i$ grows
at most linearly and its generator increments have integrable
polynomial growth. Localization and the bounded-network moment
estimates justify stationary generator testing. With
\begin{equation}
 L_\epsilon=L_0+\epsilon L_1,\qquad
 L_1=a h_2(P_2)\Delta_{M_1}-b h_1(P_1)\Delta_{M_2},
\end{equation}
put $k_i=L_1Rq_i$. Then
\begin{equation}
 \E_\epsilon q_i=\epsilon\E_\epsilon k_i,\qquad
 \E_0k_1=aU_1-bV_1,\quad
 \E_0k_2=aV_2-bU_2.
 \label{eq:poissonresponse}
\end{equation}
The two RNA birth channels acting separately on the product
Poisson solution give the last two identities. This also shows why
the Fano deviations start at quadratic order after multiplying
by the signed transcription gain.

\subsection{Weighted increment estimate}

Any function from integer counts to $[-1,1]$ has discrete Lipschitz
constant at most two. Coupling unregulated descendant families,
for $d\geq1$, gives
\begin{align}
 |\mathcal T_j(t)h_j(z)-\mathcal T_j(t)h_j(y)|
 &\leq2e^{-t}(\Delta p_j+t\Delta m_j),\nonumber\\
 |\Delta_M\mathcal T_j(t)h_j(z)|&\leq2te^{-t},\nonumber\\
 |\Delta_M\mathcal T_j(t)h_j(z)
        -\Delta_M\mathcal T_j(t)h_j(y)|
 &\leq4e^{-t}(\Delta p_j+t\Delta m_j),
\end{align}
where $\Delta m_j=|m_j-m_j'|$, $\Delta p_j=|p_j-p_j'|$
and $y$ is the primed state. For the feedforward term in $k_i$,
integrating the product difference costs at most
$g[\Delta m_j+(19/2)\Delta p_j]$.

For the return product $h_i u_i\,\Delta_M\mathcal T_jh_j$,
use the exact telescoping into the difference of $u_i$, of $h_i$,
and of the extra-RNA response, with the other factors fixed in turn.
The bound $|u_i(t,y)|\leq r_1(t)(W_i'+2)$ and the integrals
\begin{equation}
 \int r_1=2,\qquad
 \int r_1e^{-t}=\frac34,\qquad
 \int tr_1e^{-t}=\frac12,\qquad
 \int te^{-2t}=\frac14
\end{equation}
give
\begin{equation}
 g\left[\Delta m_i+
   \left(\frac92+2W_i'\right)\Delta p_i+
   (4+2W_i')\Delta m_j+(6+3W_i')\Delta p_j\right].
\end{equation}
Adding the feedforward bound therefore yields
\begin{align}
 |k_i(z)-k_i(y)|
 &\leq g\left[\Delta m_i+
   \left(\frac92+2W_i'\right)\Delta p_i+
   (5+2W_i')\Delta m_j+
   \left(\frac{31}{2}+3W_i'\right)\Delta p_j\right]\nonumber\\
 &\leq16g(1+W_i')
       (\Delta m_i+\Delta p_i+\Delta m_j+\Delta p_j).
\end{align}
Applying this to the stationary coupling
and Eq.~\eqref{eq:weightedmoment} gives
\begin{equation}
 |\E_\epsilon k_i-\E_0k_i|\leq440\epsilon g^2.
\end{equation}
Exact first-moment balance gives
$|\E_\epsilon Z_i-\bar z_i|\leq2\epsilon g$,
and Eq.~\eqref{eq:envelope} gives
$|\E_\epsilon h_j-\eta_j|\leq4\epsilon g$.
Subtracting the mean product in $\E_\epsilon q_i$, and using
Eq.~\eqref{eq:poissonresponse}, proves
\begin{equation}
 \left|\Cov_\epsilon(Z_i,h_j)-\epsilon\E_0k_i\right|
 \leq448\epsilon^2g^2 .
\end{equation}
Multiplication by $\epsilon a$ for gene one and by $-\epsilon b$
for gene two proves Eq.~\eqref{eq:remainder}. The estimate is
specific to the fixed unequal pair and is not an arbitrary-kinetics
remainder theorem.
